\section{Diagnostics, failure modes, and interpretation}
\label{ch:diagnostics}

\subsection{Three layers of verification}

A credible result needs three different checks:
\begin{description}
  \item[Accounting.] Do input margins, exposure stocks, routes, and edge flows
  describe one baseline?
  \item[Mathematics.] Does the code implement the derivative and closure
  identities?
  \item[Empirics.] Does the calibrated model describe outcomes not imposed by
  construction?
\end{description}
The package automates the first two. The third requires application-specific
evidence.

\subsection{Route diagnostics}

The run manifest records:
\begin{itemize}
  \item route spectral radius;
  \item absorption error;
  \item bilateral row and column errors;
  \item flexible-route edge reconstruction error;
  \item fixed-route edge reconstruction error.
\end{itemize}
The spectral radius must be below one. A value close to one can produce a
well-defined but poorly conditioned resolvent in which traffic circulates for
many expected steps before absorption.

Large margin or edge errors usually indicate inconsistent inputs, reversed
orientation, duplicated arcs, or an undeclared transformation. They should not
be repaired by increasing the tolerance.

\subsection{Linear-system diagnostics}

Every closure reports the condition numbers of its transport and spatial
systems. The default upper gate is \(10^{12}\). This is a permissive numerical
ceiling rather than a target. Results near the gate can be sensitive to small
changes in data or parameters.

The package also records solve residuals and finiteness checks. A low residual
with a very high condition number does not establish a stable economic
derivative. Both should be examined.

\subsection{Identity diagnostics}

For \code{decompose}, inspect:
\begin{itemize}
  \item maximum inverse-gap error;
  \item closure-ladder error;
  \item realized and primitive Hulten-accounting errors;
  \item closure-Jacobian reconstruction error;
  \item analytical channel reconstruction error;
  \item primitive transport residual.
\end{itemize}
The package acceptance target for algebraic identities is \(10^{-10}\). An
identity failure is a code or model-construction problem. It is not an
interesting economic result.

\subsection{Finite differences}

An analytical derivative should be compared with independently solved
nonlinear equilibria. For a central difference with step \(h\),
\[
 E_p^{FD}
 =-\frac{\log W(\theta_p+h)-\log W(\theta_p-h)}{2h}.
\]
The nonlinear solver must re-solve all equilibrium and transport quantities.
Reusing the analytical derivative inside the counterfactual would make the
test circular.

The package tests efficient, externality, congestion, asymmetric, multiple
mode, unique-route, and parallel-route cases. The accepted RSUE report uses a
deterministic traffic-stratified sample of road arcs. The absolute derivative
error must be below \(10^{-6}\), and nonlinear solve residuals below
\(10^{-10}\).

\subsection{First-order accuracy}

Finite-difference verification establishes the derivative at the baseline. To
assess a one-percent policy, compare:
\[
 \Delta\log W_{\mathrm{linear}}=-E_p\Delta\theta_p
\]
with a nonlinear counterfactual at \(0.1\), \(1\), and \(5\) percent where
feasible. Errors that shrink quadratically as the shock shrinks are consistent
with a correct local derivative.

The package does not automatically bless a five-percent intervention because
the one-percent derivative is accurate. The analyst must report the scale at
which the approximation is used.

\subsection{Sensitivity}

A sensitivity path should report both average effects and rank stability. A
parameter can scale all links similarly while leaving the ordering nearly
unchanged, or leave the average stable while reshuffling the top links.

The paper exercises vary \(\alpha\), \(\beta\), net dispersion, \(\eta\), and
road congestion in the headline model. Common road/terminal and terminal-only
paths are extension diagnostics. Each point must remain on the declared
regular branch and pass all closure gates.

\subsection{Common failures}

\begin{longtable}{p{0.30\textwidth}p{0.29\textwidth}p{0.33\textwidth}}
\toprule
Message or symptom & Likely cause & Correct response \\
\midrule
\endhead
Route kernel is not contractive & Cyclic continuation shares do not absorb &
Revisit flow accounting or local absorption; do not truncate paths silently. \\
Stock disagreement & Incoming and outgoing flows imply different exposure &
Document and perform a balancing transformation outside the package. \\
Policy edge lacks active mode & Zero or missing policy-mode flow & Remove the
policy or provide a defensible positive-flow baseline. \\
Missing terminal ID & Positive terminal congestion on an unidentified terminal
& Add explicit endpoint terminal identifiers. \\
Condition-number gate & Near-singular modal, transport, or spatial system &
Review parameters, baseline support, and model identification. \\
Undefined effective pass-through & Realized derivative is numerically zero &
Use the level pass-through component; do not force a ratio. \\
Physical link rejected & Missing or inconsistent reciprocal direction &
Report directed policies or correct the physical-link identifier. \\
\bottomrule
\end{longtable}

\subsection{What to report}

At minimum, an empirical application should disclose:
\begin{itemize}
  \item network and flow construction;
  \item units and every transformation;
  \item policy primitive, direction, mode, and scale;
  \item parameter sources;
  \item modal and congestion specifications;
  \item route, accounting, condition, and identity diagnostics;
  \item sensitivity of levels and ranks;
  \item external-validation evidence and its limitations.
\end{itemize}
This information is part of the economic result, not software housekeeping.
